% Chapter 9 -- Conclusion and Future Work
% Standalone, compilable draft (FPT-template article class + booktabs), sibling to
% sec_1_introduction.tex / sec_8_discussion.tex. Synthesis-only chapter: it answers the
% three research questions posed in Chapter 1 one by one, each pointing to the chapter that
% settled it; restates the three claimed contributions compactly; and turns the boundaries
% named in Chapter 8 into a concrete future-work list. No tables or figures of its own.
%
% MERGE/NUMBERING NOTE: this file opens with \section{Conclusion and Future Work} for
% standalone compilation and numbering, but at final assembly the assembler promotes it
% to a chapter (Chapter 9). Forward/back references to sibling chapters and their tables
% use typed numerals ("Chapter~7", "Section~7.6", "Table~4.5") because the chapter files
% do not share a label namespace when built separately.
%
% Honesty-phrasing rules applied (personal_agenda P1):
%   - recall denominator restated (>=10 surviving points, Sec 4.1; B2 vs-all-cars ~0.54)
%   - completion held-out outcome kept at PARTIALLY HOLDS (Finding #42), never "confirmed"
%   - moving cars = plausibility only, never "accuracy" (Finding #44)
%   - negative results (#21/#24/#25/#30/#16/#17/#19) presented as findings, not apologised for
%   - #37 as a metric-design lesson, not an erratum
%   - claims kept within the Section-6 claims map (C1-C11); nothing beyond C9's scope
%
\documentclass[a4paper,12pt]{article}
\usepackage[english]{babel}
\usepackage[utf8]{inputenc}
\usepackage[T1]{fontenc}
\usepackage{lmodern}
\usepackage{amsmath,amssymb}
\usepackage{booktabs}
\usepackage{float}
\usepackage{siunitx}
\usepackage{microtype}
\usepackage[margin=2.5cm]{geometry}

\begin{document}

\section{Conclusion and Future Work}
\label{ch:conclusion}

This thesis studied an interpretable, near-annotation-free pipeline for detecting and
completing cars in LiDAR point clouds, and asked what such a system can and cannot
establish. The detection side is a training-free clustering-plus-classifier pipeline whose
recall was root-caused rather than tuned away; the completion side is a learned model
whose value on real data had to be measured before it could be claimed, because no valid
real-data metric existed at the outset. Chapter~8 drew the boundary around each result.
This chapter states the answers in the form the questions were asked, restates what the
thesis contributes, and turns the most consequential of those boundaries into future work.

\subsection{Answers to the research questions}
\label{sec:conclusion-rqs}

Chapter~1 posed three research questions. Each is answered by a specific chapter, and each
answer is reported at the reach the evidence actually supports.

\paragraph{RQ1 --- evaluating completion on real LiDAR.}
The measurement question is answered in Chapter~7. Reference-based Chamfer distance against
an accumulated pseudo-ground-truth is \emph{invalid} on real cars: because the accumulated
reference is itself occluded on the same side as the input, the raw partial cloud scores
best on that metric, so a lower Chamfer distance rewards under-completion rather than
reconstruction (Finding~\#26). The thesis instead defines a donor-frame occluded-side
coverage metric: it completes a car from one frame's pipeline cluster and scores how much
of the surface seen only in \emph{other} frames of the same static car the completion
recovers, paired with a hallucination guard against surface invented outside the car's
amodal box. The metric passes a four-item validation gate, and the per-band form of that
guard exposed a failure the original pooled statistic was blind to --- a metric-design
lesson about band-blind gates, not an erratum to the direction it validated
(Findings~\#32/\#37). This is a validated, leakage-free way to score completion quality on
real LiDAR in this setting.

\paragraph{RQ2 --- does completion recover unseen surface, and does the benefit hold out?}
Also answered in Chapter~7, and reported exactly as it fell. On real sequence-08 cars the
donor metric shows that PCN recovers genuinely occluded surface several times above the
symmetry-mirror baseline (occluded-side coverage at \SI{0.1}{\metre} rising from
\num{0.000} for the raw partial and \num{0.043} for the mirror baseline to \num{0.304} for
the completion, per-car medians), rather than merely redistributing the points it was
given, and the completed cloud improves downstream amodal bounding boxes on static,
gate-passed cars in the compact and normal length bands (Findings~\#29/\#32). Under a
pre-registered test on the held-out sequence~00, the benefit \textsc{partially holds}
(Finding~\#42): the two primary metrics replicate on the held-out data
(BEV IoU \num{0.739}~$\rightarrow$~\num{0.766}, $p=\num{1.6e-3}$; donor coverage
\num{0.000}~$\rightarrow$~\num{0.413}), but a coverage gap prevents a full \textsc{holds} ---
the long length band was empty on sequence~00 (no car of at least \SI{4.6}{\metre} survived
the static-car construction), so long-band behaviour was untestable off the tuning sequence.
The compact-band hallucination guard also fails here, as it does on sequence~08. Moving cars
are covered by a plausibility check only, never an accuracy claim (Finding~\#44); that
plausibility scoping is stated in Section~7.5 and the held-out verdict in Section~7.6, and
neither is softened here.

\paragraph{RQ3 --- what limits recall, and is synthetic pretraining necessary?}
The first half is answered in Chapter~5, the second in Chapter~4. Recall is limited not by
the classifier but by the clustering stage: HDBSCAN splits single cars into multiple
fragments, with \SIrange{31}{37}{\percent} of ground-truth cars split (Finding~\#23), and
the classifier's role is to buy precision, which it does at a cost to recall that the track
filter then partly recovers (geometric-only recall \num{0.775}, classifier \num{0.677},
track filter \num{0.730}; Findings~\#47/\#49). A search of repair strategies --- larger
minimum cluster size, fragment merging, adaptive clustering, BEV clustering, temporal accumulation,
threshold lowering --- fails to generalise, each for a structural reason, and is reported
as a chain of negative findings (Findings~\#21/\#24). The one lever that helped was the
clustering resolution itself, which lifted recall from \num{0.699} to \num{0.730} and
showed the earlier ``hard limit'' was partly a resolution artifact (Finding~\#34); a
residual limit remains because density clustering has no notion of objectness. Recall is
stated throughout against the car-only $\geq\!10$-surviving-point denominator of
Section~4.1; against all annotated cars the effective figure is lower (approximately
\num{0.54}, Finding~\#48). On the second half, synthetic pretraining is \emph{redundant}: a
from-scratch classifier matches a synthetically pretrained one once real training clusters
are plentiful (of order \num{420000} mined clusters), and a cross-domain matrix shows the
sim-to-real gap is total and symmetric --- car $F_1$ collapses to \num{0.000} in every
off-diagonal cell (Findings~\#25/\#30). This replaces the two-stage-training contribution
promised in the proposal with a negative result.

\subsection{Contributions}
\label{sec:conclusion-contributions}

Restated compactly, the thesis makes three claimed contributions, all methodological or
diagnostic rather than architectural; it does not claim novelty for the pipeline itself.

\begin{enumerate}
  \item \textbf{A leakage-free metric for completion quality on real LiDAR.} The
  demonstration that accumulated-pseudo-ground-truth Chamfer is invalid on real cars
  (Finding~\#26), together with the donor-frame occluded-side coverage metric, its
  hallucination guard, and its four-item validation gate (Findings~\#32/\#37); and, applied
  through that metric, the evidence that completion recovers occluded surface several times
  above the mirror baseline and improves amodal boxes on static gate-passed cars, partially
  holding on a pre-registered held-out sequence (Findings~\#29/\#42). This is the
  completion half's central contribution (Chapter~7).
  \item \textbf{A root-cause account of the recall bottleneck, with a chain of negative
  repair results.} Recall is clustering-bound, not classifier-bound (Finding~\#23); a search
  of repair strategies fails to generalise, each for a structural reason, with the
  clustering resolution the only lever that helped (Findings~\#21/\#24/\#34). The negatives
  are reported as findings in their own right (Chapter~5).
  \item \textbf{Evidence that synthetic pretraining is redundant for the classifier at
  scale.} A from-scratch classifier matches a synthetically pretrained one once real
  clusters are plentiful, and the sim-to-real gap is total and symmetric
  (Findings~\#25/\#30). This directly replaces a proposal contribution (Chapter~4).
\end{enumerate}

Alongside these, the thesis carries engineering contributions that it does not claim as
scientific results --- the modular pipeline and its runtime characterisation, the amodal
ground-truth box builder, and the reproducibility infrastructure --- and a set of documented
negative-result lines (architecture swaps and real-data fine-tuning for completion,
Findings~\#16/\#17/\#19/\#28; a sparse-fallback length model, Finding~\#40).

\subsection{Future work}
\label{sec:conclusion-future}

Chapter~8 named the boundaries of the present results. Three of them are the natural next
steps, each already scoped by a finding rather than left open.

\paragraph{A real-time variant.} The pipeline runs offline at roughly
\SI{623.5}{\milli\second} per scan (\num{1.6}~frames/s; Section~4.5), dominated by the
classical CPU stages. The runtime work that was done (Findings~\#33/\#34) established that
the $\leq\SI{400}{\milli\second}$ target is not reachable through the authorised
exact-output and voxel-resolution tiers, because the post-ground object cloud is
intrinsically sparse. Reaching real time would require a learned detector or GPU
clustering, deliberately not pursued here; it is a different system and is future work, not
a claim of this thesis.

\paragraph{Real-data fine-tuning for completion.} The completion model was trained on
synthetic KITTI-like partials only, which is at once the source of the donor metric's
validity and a ceiling on real-data crispness: on the sparsest, most one-sided inputs the
completions stay flat, which Finding~\#26 attributes to the remaining synthetic-to-real
domain gap. Closing it would call for real-data fine-tuning, which is itself a documented
negative-result line in this project (Findings~\#16/\#17/\#19); it is a contingency
direction, to be attempted only with the leakage controls that make real-data supervision
safe.

\paragraph{Long-band amodal ground truth on a second held-out sequence.} The held-out
replication (Section~7.6) partially held in part because the long length band was empty on
sequence~00 --- no car of at least \SI{4.6}{\metre} survived the well-observed-static-car
construction --- so the per-car length estimate's behaviour on long cars could not be tested
off the tuning sequence. Building amodal ground truth for a third sequence with long cars
present would close that single coverage gap and upgrade the long-band claim from a
sequence-08 result to a replicated one. It is deferred here as a bounded data-construction
task, not a change to the method.

\subsection{Closing remarks}
\label{sec:conclusion-closing}

The thesis draws a deliberately narrow boundary and defends it. On the completion side it
contributes a validated, leakage-free real-data metric and pre-registered, partially
replicated evidence that learned completion adds genuinely unseen surface and improves
amodal boxes on static cars. On the detection side it contributes a root-caused account of
why a training-free clustering detector is recall-limited, a chain of negative repair
results reported as findings, and the demonstration that synthetic pretraining is redundant
at scale. None of these rests on an additional point of metric improvement; each is defended
by explanation and, where the design allowed it, by a pre-registered test whose outcome was
reportable either way. That is the standard the work set for itself, and it is the standard
on which it asks to be judged.

\end{document}
